% ===================================================================
%  Αρχείο: 1851_SOLUTION.tex (Fragment)
%  Αυτόματη μετατροπή μέσω doc2latex.py
% ===================================================================

ΘΕΜΑ 2

Σε τρίγωνο 
% TODO: Μετατροπή σχήματος σε TikZ/PGFPlots
\begin{figure}[htbp]
\centering
\includegraphics[width=0.8\linewidth]{../extracted/1851_SOLUTION/image1.pdf}
% \caption{Σχήμα 1}
\end{figure}
 η προέκταση της διχοτόμου της $\widehat{\Gamma}$ και της εξωτερικής γωνίας του $\widehat{Β}$, τέμνονται στο Ε. Δίνεται ότι 
% TODO: Μετατροπή σχήματος σε TikZ/PGFPlots
\begin{figure}[htbp]
\centering
\includegraphics[width=0.8\linewidth]{../extracted/1851_SOLUTION/image2.pdf}
% \caption{Σχήμα 2}
\end{figure}


\begin{alist}
\item Να αποδείξετε ότι το τρίγωνο ΓΒΕ είναι ισοσκελές \monades{12}

\item Να υπολογίσετε τις γωνίες του τριγώνου
% TODO: Μετατροπή σχήματος σε TikZ/PGFPlots
\begin{figure}[htbp]
\centering
\includegraphics[width=0.8\linewidth]{../extracted/1851_SOLUTION/image1.pdf}
% \caption{Σχήμα 3}
\end{figure}
 . \monades{13}


% TODO: Μετατροπή σχήματος σε TikZ/PGFPlots
\begin{figure}[htbp]
\centering
\includegraphics[width=0.8\linewidth]{../extracted/1851_SOLUTION/image3.png}
% \caption{Σχήμα 4}
\end{figure}
\end{alist}
